\chapter{Introduzione: ARM64 vs x86 — Gestione della Memoria di Processo}

\section{Layout della Memoria di Processo su Windows}

% TODO — da scrivere dopo raccolta artifact scenario 1 (Baseline)
% Coprire: stack, heap, moduli, TEB; indirizzi tipici su Win11 ARM64

\section{Calling Convention ARM64 vs x86\_64}

% Registri argomento: x0–x7 (ARM64) vs rcx/rdx/r8/r9 (x86_64)
% Caller-saved vs callee-saved
% Allineamento SP a 16 byte — obbligatorio su ARM64

\section{La Catena x29/x30}

% Prologo compilatore C: stp x29, x30, [sp, #-N]!
% Cosa finisce in memoria, come si attraversa per stack walking
% Riferimento diretto: sezione 2.1 (Baseline dump WinDbg)

\section{ARM64 vs x86: Differenze Fondamentali}

% Istruzioni a 4 byte fissi → ogni return address valido è multiplo di 4
% Niente RET da 1 byte
% LR è registro (x30), non sullo stack fino al prologo

\section{Strutture \texttt{.pdata} e \texttt{.xdata}}

% RUNTIME_FUNCTION: BeginAddress, EndAddress, UnwindData
% Unwinding e stack walking via RtlVirtualUnwind
% Return address fuori da .pdata = anomalia rilevabile
% Riferimento: sezione 2.2 (Single-Frame Spoofing)

\section{Il TEB come Guardia dello Stack}

% StackBase in x18+0x08, StackLimit in x18+0x10
% Windows usa questi per validare SP durante stack walking
% Chiave per capire lo Stack Pivoting (sezione 2.4)
